%\ifodd\value{page}\hbox{}\newpage\fi
\chapter*{Ṛgveda 10.125 — Devī Sūkta, Mantra 3}
\addcontentsline{toc}{chapter}{Devī Sūkta 10.125.3}

\begin{sloka}
{\sanskritfont
\begin{tabular}{c}
मया सो अन्नमत्ति यो विपश्यति यः प्राणिति य ईं शृणोत्युक्तम् ।\\
अमन्तवो मां त उप क्षियन्ति श्रुधि श्रुत श्रद्दिवं ते वदामि ॥३॥
\end{tabular}

\vspace{1.5em}

{\middlesize \iastfont
mayā so annam atti yo vipaśyati yaḥ prāṇiti ya īṁ śṛṇoty uktam |\\
amantavo māṁ ta upa kṣiyanti śrudhi śruta śraddhivaṁ te vadāmi ||3||}

\vspace{1em}

{\middlesize
\anvaya{%
मया सः अन्नम् अत्ति ।
यः विपश्यति ।
यः प्राणिति ।
यः ईम् उक्तम् शृणोति ।
अमन्तवः माम् ते उप क्षियन्ति ।
श्रुधि ।
श्रुत ।
श्रद्दिवम् ते वदामि ।
}
}

\middlesize \iastfont \itmean{%
«Per mezzo mio egli mangia il cibo:  
colui che vede, colui che respira, colui che ascolta la parola pronunciata.  
Coloro che non comprendono dimorano vicino a me senza riconoscermi.  
Ascolta! O ascoltante, ti dico ciò che è degno di fede.»%
}

}
\end{sloka}

\wordmeaning{%
\wordline{मया सो अन्नम् अत्ति}
{mayā so annam atti}
{ per mezzo di me (\textit{mayā}) egli mangia (\textit{atti}) il cibo (\textit{annam}); }%

\wordline{यः विपश्यति यः प्राणिति}
{yaḥ vipaśyati yaḥ prāṇiti}
{ colui che vede chiaramente (\textit{vipaśyati}), colui che respira (\textit{prāṇiti}); }%

\wordline{यः ईम् उक्तम् शृणोति}
{yaḥ īṁ uktam śṛṇoti}
{ colui che ascolta (\textit{śṛṇoti}) la parola pronunciata (\textit{uktam}); }%

\wordline{अमन्तवः माम् उप क्षियन्ति}
{amantavaḥ māṁ upa kṣiyanti}
{ coloro che sono privi di comprensione (\textit{amantavaḥ}) dimorano vicino (\textit{upa kṣiyanti}) a me senza riconoscermi; }%

\wordline{श्रुधि श्रुत श्रद्दिवम् ते वदामि}
{śrudhi śruta śraddhivam te vadāmi}
{ ascolta (\textit{śrudhi}), o ascoltante (\textit{śruta}); ti dico (\textit{vadāmi}) ciò che è degno di fede (\textit{śraddhivam}); }%
}

Nel terzo mantra la Dea esplicita la sua funzione come principio operativo della vita quotidiana. Non è solo presenza cosmica o sovranità astratta: è ciò attraverso cui si mangia, si respira, si vede e si ascolta. Ogni funzione vitale è resa possibile “per mezzo di lei”. Il verso introduce però anche una tensione: molti “dimorano vicino” alla Dea senza riconoscerla. La sua presenza è immediata e costante, ma non automaticamente conosciuta. La mancanza non è di vicinanza, ma di comprensione. Per questo il mantra si chiude con un appello diretto: “ascolta”. La Dea non impone, non costringe; invita all’ascolto consapevole. Ciò che ella comunica è detto “degno di fede”: non perché imposto da autorità esterna, ma perché riconosciuto come vero da chi ascolta con attenzione. In questo passaggio il Devī Sūkta lega in modo indissolubile vita, parola e conoscenza.

Il termine \textit{śraddhivam}, derivato da \textit{śraddhā}, indica ciò che è “degno di fede”, ma questa fede non va intesa come credenza cieca o adesione passiva. In ambito vedico, \textit{śraddhā} è una disposizione interiore che rende possibile il riconoscimento del vero: una forma di fiducia attiva, lucida, che nasce dall’ascolto e dalla risonanza interiore. Quando la Dea afferma “ti dico ciò che è \textit{śraddhivam}”, non si pone come autorità esterna che impone un contenuto, ma come principio che invita a un atto di affidamento consapevole. La verità, qui, non è qualcosa da possedere, ma qualcosa a cui ci si apre. \textit{Śraddhivam} segna dunque la soglia tra la semplice vicinanza alla presenza divina e il suo effettivo riconoscimento: senza questa fiducia vigile, si può “dimorare accanto” al sacro e tuttavia restarne ciechi. In questo senso, il mantra afferma che la conoscenza più profonda non nasce dalla forza o dalla costrizione, ma da un ascolto che accetta di fidarsi e, proprio per questo, comprende.